\documentclass[11pt]{article}

\usepackage[T1]{fontenc}
\usepackage[utf8]{inputenc}
\usepackage{lmodern}
\usepackage{amsmath,amssymb}
\usepackage[margin=0.9in]{geometry}
\usepackage{enumitem}
\usepackage[hidelinks]{hyperref}
\usepackage{xurl}

\setlist{leftmargin=*,itemsep=0.3em,topsep=0.35em}
\setlength{\parindent}{0pt}
\setlength{\parskip}{0.55em}
\setlength{\emergencystretch}{3em}

\newcommand{\loc}[1]{\textbf{Location:} #1}
\newcommand{\problem}{\textbf{Problem:}}
\newcommand{\impact}{\textbf{Why it matters:}}
\newcommand{\fix}{\textbf{Suggested fix:}}

\begin{document}

\begin{center}
  {\Large\bfseries Technical and Scientific Review}\par\vspace{0.35em}
  {\large Hybrid Quantum-Classical ADMM-QAOA for Transmission
  Reconfiguration on a Source-Derived Meshed Sub-System}\par\vspace{0.35em}
  Review date: July 24, 2026
\end{center}

\section*{Overall assessment}

The manuscript is substantially more careful than a typical small quantum
optimization case study: it separates QUBO optimality from topology and AC
validation, corrects the QUBO-to-Ising map, exposes the failure of the
historical radiality surrogate, and excludes unarchived ADMM and hardware
results.  Nevertheless, it is not yet submission-ready.  The repository's
own artifact check currently fails; the physical network model does not
consistently represent the transformer banks and generation described in the
paper; the proposed ADMM loop is not an ADMM decomposition of the hard-radial
problem as written; and several evidence and statistical statements are
stronger than the retained artifacts support.  These are major-revision
issues rather than copyediting details.

The classifications below distinguish \emph{definite errors},
\emph{unsupported claims}, \emph{likely issues}, and matters that
\emph{need external verification}.

\section*{Technical findings}

\subsection*{1. The advertised artifact lineage does not pass its own validator
(definite error; submission blocker)}

\loc{Contribution 4, the Results evidence-status subsections, Code and data
availability, and \path{koshi_admm_qaoa/results/artifact_manifest.json}.}

\problem{} The manuscript states that one strict JSON file is the source of
truth and that the manifest records the SHA-256 hashes of the derived outputs.
However, running
\begin{verbatim}
python koshi_admm_qaoa/artifact_pipeline.py --check
\end{verbatim}
fails.  The validator reports that \path{results/benchmark.json} differs from
the manifest; all generated tables, macros, the summary, and the copied study
protocol differ from their recorded hashes; and several audited code files,
including \path{qubo_builder.py}, \path{artifact_pipeline.py},
\path{study_protocol.py}, and \path{statistics_report.py}, differ from the
hashes frozen in the benchmark artifact.  The artifact also records a dirty
working tree at migration.  The local unit suite passes 27 tests, but those
tests do not repair the broken evidence lineage.

\impact{} The headline numbers may be internally plausible, but the current
repository cannot demonstrate that the compiled manuscript, generated files,
frozen JSON, and audited code belong to one reproducible state.  This directly
contradicts a central contribution and prevents an auditable submission.

\fix{} Freeze a clean commit, regenerate the authoritative JSON and every
derived file from that commit, rewrite the manifest last, and require both
\path{artifact_pipeline.py --check} and \path{study_protocol.py --check} to
pass in a clean checkout.  Record the successful validation output in the
release or continuous-integration log.

\subsection*{2. Transformer-bank aggregation and the network summary are
internally inconsistent (definite error)}

\loc{Section ``The Eastern Nepal sub-system,'' Table~1, the SOC model, and
\path{koshi_admm_qaoa/network_data.py}.}

\problem{} The text describes the Inaruwa 400/220-kV connection as a
$3\times315$-MVA bank, but branch 0 is modeled as one 315-MVA branch with
$x=0.12(100/315)=0.0381$ pu.  If all three identical units are in parallel,
the equivalent should instead have a 945-MVA rating and approximately
$x=0.0381/3=0.0127$ pu on the common base.  The same inconsistency occurs for
the stated $2\times160$-MVA Inaruwa 220/132-kV bank: branch 1 uses 160 MVA and
$x=0.075$ pu rather than the parallel equivalent 320 MVA and $x=0.0375$ pu.
By contrast, the $2\times15$-MVA Dhungesanghu transformer is aggregated as
30 MVA, so the treatment is not even uniform within the model.

The descriptive counts also conflict with the implemented data.  The prose
says there are four transformer banks, while the model has two branches of
kind \texttt{xfmr} plus one \texttt{tie+xfmr}.  Table~1 attributes the full
173 MW of local generation to ``Sanima MT + Kabeli IPPs,'' but the code obtains
173 MW from 60 MW at Tumlingtar, 73 MW at Dhungesanghu, and 40 MW at Amarpur.

\impact{} Transformer impedance and rating directly affect voltage drops,
thermal feasibility, losses, load shedding, and the topology selected by the
continuous subproblem.  The generation-label mismatch obscures the operating
point being studied.

\fix{} State whether each physical unit is represented separately or as a
parallel equivalent and make the rating/reactance conversion consistent.
Add a generated branch-parameter/provenance table containing unit multiplicity,
length, conductor assumption, $r$, $x$, rating, tap, and source/estimate status.
Correct the transformer count and list the three generation injections
explicitly in Table~1.

\subsection*{3. The parameter provenance is overstated (unsupported claim)}

\loc{Abstract, Contributions 1 and 4, the network-data subsection, and the
claim that branch-specific limits are published ratings.}

\problem{} The manuscript says that impedance data are derived from the NEA
Year Book and describes the limits as published branch ratings.  The released
model instead declares generic ``typical conductor'' values, uses one common
$R/X$ assumption per voltage/conductor class, assigns 600 MVA to the 220-kV
circuits and 160 MVA to most 132-kV branches, and contains several line-length
comments without a source tag.  The caveat that ``several electrical
parameters'' are estimates is too broad to identify which numerical inputs
are measured, published, inferred, or assumed.

\impact{} The topology may be source-derived, but the physical feasibility and
loss calculations depend on parameters whose provenance is not presently
traceable at branch level.  This weakens both the engineering interpretation
and the paper's reproducibility claim.

\fix{} Replace ``published'' with ``assumed'' wherever a branch-specific source
cannot be documented.  Generate the proposed provenance table from the network
data and attach a citation or explicit engineering-assumption label to every
length, impedance, rating, transformer multiplicity, and injection.

\subsection*{4. The ADMM loop is not a decomposition of the stated
hard-radial optimization problem (definite conceptual mismatch)}

\loc{Eqs.~(8) and (9), Figure~2, and Sections ``Consensus ADMM decomposition''
and ``Flow-based connected-radial formulation and projection.''}

\problem{} The stated reconfiguration problem minimizes
$\Phi_{\rm SOC}(\boldsymbol z)$ subject to a connected, radial topology.  An
exact consensus split would put the indicator of that feasible set in the
$\boldsymbol z$ block.  Instead, Eq.~(9) introduces a finite heuristic prior
$g_Q(\boldsymbol z)$, and the actual ADMM configuration sets
$\lambda_{\rm card}=0$ while retaining only soft cycle, anti-islanding, and
loss biases.  The final Kruskal projection is performed after termination and
can change the binary solution.  Therefore the loop optimizes a different
nonconvex surrogate and the projected output is not an ADMM solution of
Eq.~(8).  Approximate QAOA $\boldsymbol z$ updates also fall outside the usual
exact-subproblem ADMM convergence argument.

\impact{} Even if residuals become small, this does not establish feasibility,
stationarity, or optimality for the connected-radial reconfiguration problem.
The current paper prudently withholds numerical ADMM claims, but its method
description still implies a stronger mathematical relationship than exists.

\fix{} Either call the method an ``ADMM-inspired consensus heuristic with
post-hoc spanning-tree projection,'' or put an exact spanning-tree indicator
or a provably exact penalty in the discrete block.  State explicitly what
problem the iterates solve, what convergence guarantee is absent, and report
the raw-to-projected objective and topology changes.  Figure~2 should also
show the maximum-iteration exit and use ``after termination,'' not ``after
convergence.''

\subsection*{5. Radial operation and the series-only AC model need a
transmission-specific justification (unsupported claim / likely issue)}

\loc{Introduction, Eq.~(8c), the SOC and nonlinear-recovery subsections, and
Limitations.}

\problem{} The paper begins from optimal \emph{transmission} switching but then
requires the 400/220/132-kV sub-system to become a spanning tree.  Radiality is
standard in the cited distribution-reconfiguration setting, but it is not a
generic requirement of OTS or post-contingency transmission operation.  The
manuscript gives no protection, restoration, islanding, or operating-policy
rationale for deliberately removing all mesh redundancy.

The physical checks also use a series-only network: line charging, bus shunts,
transformer phase shifts, and magnetizing branches are omitted.  Thus
``nonlinear-AC validated'' means internally consistent only for this simplified
series-only model, not validation against a full transmission $\pi$-model.

\impact{} The hard-radial constraint may define an artificial engineering task,
and omitted reactive elements can materially change voltage and MVAr behavior
at 132--400 kV.  These choices affect the paper's central claim that the case
study is a realistic transmission application.

\fix{} Provide a system-specific reason for radial emergency operation and
state the switching/security assumptions it entails; otherwise reformulate the
study as meshed OTS with connectivity but without a tree constraint.  Add line
charging and relevant shunts/transformer details, or consistently rename the
check ``series-only nonlinear AC recovery'' and add the omission to the main
limitations and abstract caveat.

\subsection*{6. The historical QUBO is not fully specified in the manuscript
(definite implementation ambiguity)}

\loc{Eq.~(13) and the three preceding QUBO-component paragraphs.}

\problem{} The paper decomposes the linear coefficient into consensus, loss,
cardinality, and anti-islanding components, but it never defines the consensus
or loss coefficients.  The implementation uses
\[
 q^{\rm cons}_\ell=\frac{\rho}{2}
 \bigl[1-2(\alpha_\ell+u_\ell)\bigr],\qquad
 q^{\rm loss}_\ell=\eta
 \frac{r_{\pi(\ell)}}{\max_m r_{\pi(m)}+10^{-9}}.
\]
The second coefficient is positive in a minimization QUBO: with fixed
cardinality it favors lower-resistance closed branches, but when
$\lambda_{\rm card}=0$ it is also a direct bias toward opening switches.  That
effect is not explained.

Two local mathematical statements are also too strong.  A finite two-edge
cycle penalty does not ``exclude'' the all-closed state; it only penalizes it
unless a sufficient penalty bound is proved.  In addition, the truncated
anti-islanding expression is exact up to a constant for degree-one as well as
degree-two buses; the failure begins at degree three.

\impact{} A reader cannot reconstruct the QUBO, interpret the loss bias, or
assess whether the soft penalties dominate competing terms from the equations
alone.

\fix{} Display every coefficient, normalization, loop set, and sign convention;
archive the actual linear vector and upper-triangular QUBO matrix for each
instance.  Replace ``excludes'' with ``penalizes'' unless a sufficient penalty
bound is supplied, and correct the degree-one/degree-two statement.

\subsection*{7. The frozen artifact does not support the phrase ``the
exact-QUBO solver's returned bitstring'' (definite evidence mismatch)}

\loc{Abstract, the legacy SOC-audit Results subsection, Lesson 2, and
\path{koshi_admm_qaoa/results/benchmark.json}.}

\problem{} For every retained size, the fields
\path{one_minimizer_bits_variable_order} and \path{returned_bits} are null.
The artifact marks connectivity as ``inferred from legacy loss,'' and Table~A.6
explicitly defines D as either disconnected or a non-finite record whose cause
was not separately archived.  The manuscript nevertheless states that the
topology induced by the solver's returned minimizer is disconnected at
$n=6,10,12$.

Current exhaustive reconstruction does find disconnected global minimizers at
those sizes, so the qualitative topology conclusion may be correct.  It is not,
however, evidence about the particular historical bitstring returned by the
legacy solver, especially while the audited QUBO code hash no longer matches
the frozen artifact.

\impact{} The wording crosses the paper's own boundary between retained
evidence and reconstructed audit evidence.

\fix{} Regenerate and archive the complete minimizer set, the exact solver's
selected bitstring, variable order, and topology diagnostics.  Until then,
attribute the statement to a separately versioned current enumeration and do
not describe it as the retained solver output.

\subsection*{8. The prospective confidence-interval description does not
match the code, and seed labels do not by themselves create pairing
(definite error)}

\loc{Prospective protocol, Table~3, \path{statistics_report.py}, and
\path{post_contingency_pipeline.py}.}

\problem{} The manuscript and generated-table caption call the median-gap
interval a ``paired-seed bootstrap'' interval.  The function
\path{summarize_trials} actually performs an ordinary within-method percentile
bootstrap of the median.  A separate function computes seed-wise differences,
but its result is not the interval placed in the table.  Moreover, assigning
the same integer seed to QAOA, QRAO, and SA does not automatically define a
meaningful matched experimental unit: the algorithms consume different random
streams and do not share a common stochastic disturbance.

\impact{} The prospective table would mislabel its uncertainty calculation,
and a paired interpretation could produce unjustified precision for
between-method comparisons.

\fix{} Call the displayed intervals unpaired percentile-bootstrap intervals.
If pairwise method contrasts are desired, predeclare the contrast, justify a
real blocking or common-random-number design, and report the bootstrap interval
for the seed-wise difference rather than the separate method medians.

\section*{Meaningful editorial and presentation findings}

\subsection*{Algorithm-flow wording}

\loc{Figure~2 caption and the ADMM iteration contract.}

\problem{} The caption says projection occurs ``after convergence,'' while the
implementation also projects after the 30-iteration limit.  The diagram has no
maximum-iteration exit.

\impact{} A nonconverged iterate can be mistaken for a converged ADMM result.

\fix{} Use ``after termination,'' add the iteration-limit branch, and require
the termination reason beside every final result.

\subsection*{Reference integrity}

\loc{References 11, 14, and the compiled bibliography generally.}

\problem{} The Bhandari entry has incorrect or unverified initials and omits
the page range; the journal metadata gives Ganesh Bhandari, Bishal Rimal, and
Sandeep Neupane, pp.~75--80, DOI \texttt{10.3126/tj.v2i1.32842}.  The Morstyn
entry leaves a ``Verify volume/pages'' note in the publication bibliography
and omits DOI \texttt{10.1109/TSG.2022.3200590}, although volume 14(2),
pp.~1093--1102 is verifiable.  The Bhandari entry similarly prints ``Author
first names to be confirmed from source.''

\impact{} These are visible submission defects and weaken citation integrity.

\fix{} Correct the metadata, add the missing DOI/page range, and remove all
author-facing verification notes before submission.

\section*{Proofing sweep}

The bundled mechanical scan was run on the compiled PDF, the manuscript
source, and the bibliography, and the reported hits were spot-checked.  The
apparent doubled periods around paragraph headings were PDF-extraction false
positives.  The following high-confidence defects remain:

\begin{itemize}
  \item \textbf{Compiled references:} BibTeX case conversion produces
  ``aC-feasible,'' ``iEEE,'' ``sOC,'' lower-case ``qiskit,'' lower-case
  ``nepal,'' and ``3-admm-h.''  Protect technical and proper names with braces,
  for example \texttt{\{AC\}}, \texttt{\{IEEE\}}, \texttt{\{SOC\}},
  \texttt{\{Qiskit\}}, \texttt{\{Nepal\}}, and \texttt{\{3-ADMM-H\}}.
  \item \textbf{Reference notes:} Remove internal instructions such as
  ``Verify volume/pages against final record'' and ``Author first names to be
  confirmed from source.''
  \item \textbf{Reference titles and names:} Protect geographic proper names
  such as Hetauda--Dhalkebar--Inaruwa and Integrated Nepal Power System so the
  bibliography style does not lowercase them.
\end{itemize}

\section*{Checks performed without a concrete error}

The required notation/terminology drift audit was performed.  The core roles
of $\boldsymbol\alpha$ (continuous ADMM copy), $\boldsymbol z$ (binary copy),
$\boldsymbol b^{\rm raw}$, $\boldsymbol b^{\rm proj}$, the map $\pi$, and the
run-level versus shot-level hit statistics remain verbally consistent across
the manuscript.  No initial/final or source/observer descriptor conflict was
found for these symbols.

The QUBO-to-Ising coefficients in Eq.~(15), including the additive constant,
are algebraically correct for $z_\ell=(1-Z_\ell)/2$.  The branch-flow angle
formula uses \texttt{atan2} with the correct numerator and denominator for the
stated real-tap series convention, so no quadrant ambiguity remains.  The
single-commodity-flow balance, gate bounds, and edge-count argument correctly
characterize a spanning tree when the mandatory subgraph is a forest.  The
SOC cone, sign conventions, and per-unit dimensions were also checked; no
additional concrete algebraic or dimensional error was identified beyond the
physical-parameter issues above.

\section*{Highest-priority fixes}

\begin{enumerate}
  \item Restore a clean, validator-passing artifact/code/manuscript lineage.
  \item Correct transformer aggregation and publish branch-level parameter
  provenance before interpreting any physical result.
  \item Reframe the ADMM method as a heuristic or make its discrete block
  enforce the stated connected-radial feasible set exactly.
  \item Justify radial transmission operation and qualify or extend the
  series-only AC model.
  \item Archive exact bitstrings/QUBO matrices and correct the prospective
  bootstrap description.
  \item Repair the bibliography metadata and capitalization.
\end{enumerate}

\section*{Items needing external verification}

\begin{itemize}
  \item The novelty statement that no prior work targets a real, meshed,
  multi-voltage-class transmission sub-system with provenance-traceable data
  and nonlinear validation requires a documented literature search; the cited
  examples alone do not prove absence of prior work.
  \item The citation to Fisher \emph{et al.} does not obviously establish the
  specific claim that Eq.~(8) is NP-hard because it contains a
  minimum-loss radial network-design problem as a special case.  Supply a
  direct complexity reference or a short reduction, and distinguish general
  OTS from the added spanning-tree restriction.
\end{itemize}

\end{document}